\documentclass[11pt]{article}
\usepackage[a4paper,margin=25mm]{geometry}
\usepackage{amsmath,amssymb,mathtools}
\usepackage{fontspec,unicode-math}
\setmainfont{Latin Modern Roman}
\setsansfont{Latin Modern Sans}
\setmonofont{Latin Modern Mono}
\setmathfont{Latin Modern Math}
\usepackage{xurl,hyperref}
\hypersetup{colorlinks=true,urlcolor=blue,linkcolor=blue}
\newcommand{\sref}[2]{\hyperlink{source:#1:#2}{[S#2]}}
\newcommand{\cataloguescope}[1]{\paragraph{Mathematical question.}#1}
\begin{document}
% UnsolvedMath ID 8700067; source catalogue code AMR-086-0067
\setcounter{section}{8700066}
\section{Conjecture 5.4 — Loxton and van der Poorten}\label{8700067}
{\small\sffamily UnsolvedMath ID 8700067 · Graph Theory}
\subsection{Definitions and mathematical statement}

\paragraph{Shared notation.}$\mathbb N=\{1,2,\ldots\}$, and $\mathbb Z,\mathbb Q,\mathbb R,\mathbb C$ denote the integers, rationals, reals and complex numbers. Any other conventions are specified locally.


Let $0<n_0<n_1<\cdots$ be integers, $p$ a prime, and let $\mathbb F_p$ be the field of $p$ elements. Set $f(z)=\sum_{i\ge0}z^{n_i}\in\mathbb F_p[[z]]$ and $\xi=\sum_{i\ge0}10^{-n_i}\in\mathbb R$. Algebraicity of $f$ over $\mathbb F_p(z)$ means that a nonzero polynomial $P(z,T)\in\mathbb F_p[z,T]$ satisfies $P(z,f(z))=0$; irrationality means $f\notin\mathbb F_p(z)$.
\paragraph{Statement recorded in the supplied source.}
If for some prime $p$ the series $f$ is algebraic but irrational over $\mathbb F_p(z)$, must $\xi$ be transcendental over $\mathbb Q$? \sref{8700067}{2}

\subsection{Short English statement}
Does an algebraic but nonrational generating series in some finite characteristic force the corresponding decimal number with digits at those positions to be transcendental?
\subsection{Sources}
\begin{description}
\item[\hypertarget{source:8700067:1}{S1}] ULAM.AI and the UnsolvedMath Contributors, UnsolvedMath: A Curated Collection of Open Mathematics Problems, record 8700067 (AMR-086-0067). CC BY 4.0. \url{https://huggingface.co/datasets/ulamai/UnsolvedMath}.
\item[\hypertarget{source:8700067:2}{S2}] Michel Waldschmidt, \emph{Open Diophantine Problems}, arXiv:math/0312440v2 (24 January 2004), Conjecture 5.4, p.48. \url{https://arxiv.org/abs/math/0312440v2}.
\end{description}
\label{end:8700067}
\end{document}
